\documentclass[12pt,spanish]{article}
\usepackage[utf8]{inputenc}
\usepackage{babel}
\usepackage{amsmath}
\usepackage{listings}
\usepackage{mathpazo}
\usepackage{enumitem}
\usepackage{courier}
\usepackage{textcomp}
\usepackage{xcolor}
\usepackage{parskip}
\usepackage{fullpage}

\newcommand{\onelinerule}{\rule[2.3ex]{0pt}{0pt}}
\newcommand{\twolinerule}{\rule[6.2ex]{0pt}{0pt}}
\newcommand{\respuesta}{\framebox[\textwidth]{\twolinerule}}
\newcommand{\nombre}{%
  \begin{tikzpicture}[xscale=.4,yscale=.7]
    \draw (0, 0) rectangle (22, 1);
  \end{tikzpicture}%
}
%\newcommand{\rol}   {\framebox[0.3\textwidth]{\onelinerule}}
\newcommand{\rol}{%
  \begin{tikzpicture}[xscale=.4,yscale=.7]
    \draw[gray!40] ( 0, 0) grid      ( 9, 1);
    \draw          ( 0, 0) rectangle ( 9, 1);
    \draw          (10, 0) rectangle (11, 1);
    \draw (9 + .2, .5) -- (10 - .2, .5);
  \end{tikzpicture}%
}
\newcommand{\li}{\lstinline}
\providecommand{\pond}[1]{[{\small\textbf{#1\%}}]}

\lstdefinelanguage{py}{%
  classoffset=0,%
    morekeywords={%
      False,class,finally,is,return,None,continue,for,lambda,try,%
      True,def,from,nonlocal,while,and,del,global,not,with,print,%
      as,elif,if,or,yield,assert,else,import,pass,break,except,in,raise},%
    keywordstyle=\color{black!80}\bfseries,%
  classoffset=1,
    morekeywords={int,float,str,abs,len,raw_input,exit,range,min,max,%
      set,dict,tuple,list,bool,complex,round,sum,all,any,zip,map,filter,%
      sorted,reversed,dir,file,frozenset,open,%
      array,zeros,ones,arange,linspace,eye,diag,dot},
    keywordstyle=\color{black!50}\bfseries,%
  classoffset=0,%
  sensitive=true,%
  morecomment=[l]\#,%
  morestring=[b]',%
  morestring=[b]",%
  stringstyle=\em,%
}

\lstdefinelanguage{testcase}{%
  moredelim=[is][\bfseries]{`}{`},%
  backgroundcolor=\color{gray!20},%
}

\lstdefinelanguage{file}{%
  frame=single,%
}

\lstset{language=py}
\lstset{basicstyle=\ttfamily}
\lstset{columns=fixed}
\lstset{upquote=true}
\lstset{showstringspaces=false}
\lstset{rangeprefix=\#\ }
\lstset{includerangemarker=false}

\newlist{certamen}{enumerate}{1}
\setlist[certamen]{%
  label=\arabic*.,
  font=\LARGE\bfseries,%
  labelindent=-.5in,%
  leftmargin=0pt,%
  labelsep=1em%
}



\lstset{language=py}

\begin{document}
  \thispagestyle{empty}
  \section*{Miércoles 17 de octubre}

  Es posible que usted haya notado
  que la selección chilena de fútbol
  está participando de la eliminatoria
  para el próximo mundial.
  Hoy vamos a escribir código
  para analizar algunas estadísticas de esta competencia.
  (No es necesario que le guste o que sepa algo de fútbol;
  la programación nos entrega ya suficiente entretención).

  La eliminatoria es un torneo en el que participan
  nueve equipos, cada uno representando a un país sudamericano.
  En cada partido juegan dos equipos.
  El equipo que juega en su propio país se llama el \emph{local}
  y el otro se llama la \emph{visita}.
  Siempre se nombra primero al equipo local;
  por ejemplo,
  el partido Chile--Argentina es el que se jugó en Santiago,
  no el que se jugó en Buenos Aires.

  Para representar los datos de un torneo como éste,
  usaremos un diccionario cuyas llaves son los partidos.
  Cada partido es una tupla con los nombres de ambos equipos;
  por ejemplo, \li!('Chile', 'Argentina')!.

  El valor asociado a cada llave es una tupla que tiene dos valores:
  la fecha del partido y su resultado.
  La fecha, a su vez, es una tupla \li!(anno, mes, dia)!,
  y el resultado es una tupla \li!(gl, gv)!
  con los goles que anotó cada equipo.
  Si el partido aún no se juega,
  el resultado aparece reemplazado por el valor \li!None!.

  Como ejemplo, veamos los datos del partido de anoche:
  \lstinputlisting[linerange=PARTIDO\ ANOCHE-FIN\ PARTIDO\ ANOCHE]{casos.txt}

  Veamos también un partido del próximo año:
  \lstinputlisting[linerange=PARTIDO\ FUTURO-FIN\ PARTIDO\ FUTURO]{casos.txt}

  El diccionario de partidos está almacenado
  en la variable \li!partidos!,
  en un módulo llamado \li!eliminatoria.py!.

  Su equipo debe desarrollar un módulo llamado \li!torneos.py!,
  que contenga las funciones requeridas a continuación.
  Todas las funciones reciben como primer parámetro
  el diccionario de partidos.

  Los ejemplos suponen que el diccionario \li!partidos!
  ya ha sido importado del módulo \li!eliminatoria.py!.
  \newpage

  \begin{enumerate}[leftmargin=0pt]

    \item
      La función \li!contar_partidos(partidos)!
      debe retornar la cantidad total de partidos
      que hay en el torneo.
      \lstinputlisting[linerange=CONTAR-FIN\ CONTAR]{casos.txt}

    \item
      La función \li!obtener_equipos(partidos)!
      debe retornar una lista con todos los equipos
      que participan en el torneo,
      sin repetir y en orden alfabético.
      \lstinputlisting[linerange=EQUIPOS-FIN\ EQUIPOS]{casos.txt}

    \item
      La función \li!obtener_fechas(partidos)!
      debe retornar una lista con todas las fechas
      en que hay partidos,
      sin repetir y en orden cronológico.
      \lstinputlisting[linerange=FECHAS-FIN\ FECHAS]{casos.txt}

    \item
      El equipo que gana un partido recibe tres puntos,
      y el que pierde no recibe ninguno.
      Si hay un empate,
      ambos reciben un punto.
      La función \li!calcular_puntos(partidos, equipo)!
      debe retornar la cantidad total de puntos
      obtenidos por el \li!equipo! en el torneo.
      \lstinputlisting[linerange=PUNTOS-FIN\ PUNTOS]{casos.txt}

    \item
      La \emph{diferencia de goles} de un equipo
      es el total de goles que anotó
      menos el total de los que le anotaron.
      La función \li!calcular_diferencia(partidos, equipo)!
      debe retornar la diferencia de goles
      del \li!equipo! en el torneo.
      \lstinputlisting[linerange=DIFERENCIA-FIN\ DIFERENCIA]{casos.txt}

    \item
      Para decidir qué equipos clasifican al mundial,
      ellos son ordenados por puntaje.
      En caso de que dos equipos tengan el mismo puntaje,
      el que tiene la mejor diferencia de goles tiene preferencia.

      La función \li!ordenar_equipos(partidos)!
      debe entregar la lista de equipos ordenada
      según este criterio.
      \lstinputlisting[linerange=ORDENAR-FIN\ ORDENAR]{casos.txt}

      \textbf{Palta:}
      para ordenar los equipos por puntaje y luego por diferencia de gol,
      podemos aprovechar que las tuplas son ordenadas siempre
      en orden lexicográfico.
      Los pasos para lograr este ordenamiento son los siguientes:
      \begin{enumerate}
        \item
          crear una lista de tuplas \li!(puntaje, diferencia, equipo)!,
        \item
          ordenar esta lista,
        \item
          sacar ordenadamente el último elemento de cada tupla,
          e ir metiéndolos en una nueva lista.
      \end{enumerate}
      Como la primera lista tiene tuplas,
      éstas serán ordenadas primero por puntaje
      y luego por diferencia de gol.
      Por lo tanto,
      al extraer los equipos en el mismo orden,
      tendremos lo que estábamos buscando.
  \end{enumerate}

\end{document}

